\chapter{Modules and Project Organization}
\label{ch:modules}

\begin{quote}
\textit{``The purpose of abstraction is not to be vague, but to create a new semantic level in which one can be absolutely precise.''}\\
---Edsger W. Dijkstra
\end{quote}

\vspace{1em}

Novus provides a module system for organizing code into logical units with explicit visibility control. This chapter covers imports, visibility modifiers, project structure, and the standard library organization.

\section{Module Basics}
\label{sec:module-basics}

In Novus, each source file (\texttt{.novus}) is a module. The module's namespace is determined by its location in the directory structure.

\subsection{Module Naming}
\label{subsec:module-naming}

Module paths use \texttt{::} as the separator:

\begin{center}
\begin{tabular}{ll}
\textbf{File Path} & \textbf{Module Path} \\
\hline
\texttt{src/main.novus} & \texttt{main} \\
\texttt{src/renderer.novus} & \texttt{renderer} \\
\texttt{src/graphics/bitmap.novus} & \texttt{graphics::bitmap} \\
\texttt{src/sound/player.novus} & \texttt{sound::player} \\
\end{tabular}
\end{center}

The standard library uses the \texttt{std} prefix:

\begin{center}
\begin{tabular}{ll}
\textbf{File Path} & \textbf{Module Path} \\
\hline
\texttt{std/core.novus} & \texttt{std::core} \\
\texttt{std/collections/vec.novus} & \texttt{std::collections::vec} \\
\texttt{std/ffi/exec.novus} & \texttt{std::ffi::exec} \\
\end{tabular}
\end{center}

\section{Imports}
\label{sec:imports}

Use the \texttt{from ... import ...} syntax to bring items from other modules into scope.

\subsection{Basic Imports}
\label{subsec:basic-imports}

Import specific items by name:

\begin{lstlisting}
// Import specific types
from std::core import Option, Result

// Import functions and constants
from std::ffi::exec import AllocMem, FreeMem
from std::ffi::amiga_consts import MEMF_FAST, MEMF_CLEAR
\end{lstlisting}

\begin{comingfromc}
C's \texttt{\#include} literally copies header files into your source, leading to:
\begin{itemize}
    \item Include order dependencies
    \item Include guards (\texttt{\#ifndef}/\texttt{\#define}/\texttt{\#endif})
    \item Slow compilation from re-parsing headers
    \item Polluted namespaces with everything in headers
\end{itemize}

\begin{lstlisting}[language=C]
// C: Headers are textually included
#include <exec/memory.h>
#include <proto/exec.h>     // Order matters!
\end{lstlisting}

Novus imports are semantic, not textual:

\begin{lstlisting}
// Novus: Import specific symbols
from std::ffi::exec import AllocMem, FreeMem
from std::ffi::amiga_consts import MEMF_FAST
\end{lstlisting}

Benefits:
\begin{itemize}
    \item No include guards needed
    \item Order doesn't matter
    \item Only imported names are in scope
    \item Faster compilation (no re-parsing)
    \item No ``header file'' vs ``source file'' distinction
\end{itemize}
\end{comingfromc}

\subsection{Wildcard Imports}
\label{subsec:wildcard-imports}

Import all public items from a module:

\begin{lstlisting}
// Import everything from amiga_structs
from std::ffi::amiga_structs import *

// Now Task, MsgPort, Library, etc. are all in scope
\end{lstlisting}

Use wildcards sparingly---they can make it unclear where symbols come from.

\subsection{Pattern Imports}
\label{subsec:pattern-imports}

Import multiple items matching a prefix or suffix pattern:

\begin{lstlisting}
// Import all IDCMP_* constants (prefix pattern)
from std::ffi::amiga_consts import IDCMP_*

// This imports IDCMP_CLOSEWINDOW, IDCMP_MOUSEBUTTONS, etc.

// Import all *Error types (suffix pattern)
from std::error::errors import *Error

// This imports DosError, ExecError, GraphicsError, etc.
\end{lstlisting}

\subsection{Import Aliases}
\label{subsec:import-aliases}

Rename imports to avoid conflicts or for convenience:

\begin{lstlisting}
from std::collections::vec import Vec as DynamicArray
from std::collections::hashmap import HashMap as Map
\end{lstlisting}

Using the aliases:

\begin{lstlisting}
// ... (with DynamicArray and Map imported above)
var items = DynamicArray::<i32>::new()
var lookup = Map::<Str, i32>::new()
\end{lstlisting}

\subsection{Multiple Imports}
\label{subsec:multiple-imports}

Import from multiple modules at the top of your file:

\begin{lstlisting}
// Recommended import order for stdlib files:
// 1. std::core - fundamental types
from std::core import Option, Result, Drop, Clone

// 2. std::collections - data structures
from std::collections::vec import Vec

// 3. std::memory - memory management
from std::memory::block import MemoryBlock

// 4. std::strings - string types
from std::strings::core import Str

// 5. std::error - error types
from std::error::errors import ExecError

// 6. std::ffi - AmigaOS bindings
from std::ffi::amiga_structs import Task, MsgPort
from std::ffi::amiga_consts import MEMF_PUBLIC
from std::ffi::exec import AllocMem, FreeMem
\end{lstlisting}

\section{Visibility Modifiers}
\label{sec:module-visibility}

Novus has four visibility levels that control where items can be accessed.

\subsection{Private (Default)}
\label{subsec:private}

Items without a visibility modifier are private to their module:

\begin{lstlisting}
// Only visible in this file
const BUFFER_SIZE: u32 = 1024

fn helper() -> i32 {
    42
}

struct InternalData {
    value: i32,
}
\end{lstlisting}

Private items cannot be imported by other modules.

\subsection{Public (\texttt{pub})}
\label{subsec:public}

The \texttt{pub} keyword exports items, making them visible to other modules:

\begin{lstlisting}
// Exported - can be imported by other modules
pub const VERSION: u32 = 1

pub fn create_bitmap(w: u32, h: u32) -> Result<Bitmap, GraphicsError> {
    // ...
}

pub struct Bitmap {
    width: u32,
    height: u32,
    data: *u8,
}
\end{lstlisting}

Use \texttt{pub} for your module's public API---functions, types, and constants that other code should use.

\begin{comingfromc}
C uses \texttt{static} for file-local symbols and relies on headers for ``public'' APIs:

\begin{lstlisting}[language=C]
// C: static = private to this file
static int helper(void) { return 42; }

// C: non-static = visible (if declared in header)
int public_api(void) { return helper(); }
\end{lstlisting}

Novus inverts this: private by default, explicit \texttt{pub} for public:

\begin{lstlisting}
// Novus: private by default
fn helper() -> i32 { return 42 }

// Novus: pub = exported
pub fn public_api() -> i32 { return helper() }
\end{lstlisting}

The mapping:
\begin{itemize}
    \item C \texttt{static} $\approx$ Novus (default, no modifier)
    \item C non-static + header $\approx$ Novus \texttt{pub}
    \item C \texttt{extern} $\approx$ Novus \texttt{extern} (for FFI only)
\end{itemize}
\end{comingfromc}

\subsection{Internal (\texttt{internal})}
\label{subsec:internal}

The \texttt{internal} modifier makes items visible within the same package (workspace) but not exported to external code:

\begin{lstlisting}
// ... (internal visibility - illustrative)
internal const SHARED_BUFFER_SIZE: u32 = 4096

internal fn shared_helper() -> i32 {
    return 100
}
\end{lstlisting}

Use \texttt{internal} for implementation details shared between modules in a library.

\subsection{Extern (\texttt{extern})}
\label{subsec:extern}

The \texttt{extern} keyword declares items defined elsewhere, resolved at link time:

\begin{lstlisting}
// ... (extern declarations - illustrative)
extern var SysBase: *ExecBase

// Hardware register at specific address
extern var CUSTOM: *Custom at 0xdff000

// Function from another object file or assembly
extern fn asm_memcpy(dest: *u8, src: *u8, len: u32)
\end{lstlisting}

The \texttt{at ADDRESS} syntax is used for memory-mapped hardware registers. Common examples include \texttt{CUSTOM} at 0xdff000 (custom chips), \texttt{CIA\_A} at 0xbfe001, and \texttt{CIA\_B} at 0xbfd000.

\textbf{Note:} AmigaOS library functions are imported from the standard library FFI modules, not declared as \texttt{extern} in user code:

\begin{lstlisting}
// Import AmigaOS functions from std::ffi
from std::ffi::exec import AllocMem, FreeMem, Wait, Signal
from std::ffi::dos import Open, Close, Read, Write
\end{lstlisting}

External declarations have no implementation in Novus---they're provided by AmigaOS, assembly files, or C libraries.

\section{Re-exports}
\label{sec:reexports}

Re-exports allow modules to expose items from other modules as part of their own API. This is useful for creating convenient public interfaces.

\subsection{The \texttt{pub use} Syntax}
\label{subsec:pub-use}

\begin{lstlisting}
// ... (example from std/prelude.novus - illustrative)

// Re-export core types
pub use std::core::Option, Result, Drop, From

// Re-export string types
pub use std::strings::core::Str

// Re-export error types
pub use std::error::errors::DosError, ExecError, IntuitionError
\end{lstlisting}

Now users can import from \texttt{std::prelude} instead of multiple modules:

\begin{lstlisting}
// Instead of:
from std::core import Option, Result
from std::strings::core import Str
from std::error::errors import DosError

// Users can write:
from std::prelude import Option, Result, Str, DosError
\end{lstlisting}

\subsection{Multiple Re-exports}
\label{subsec:multiple-reexports}

Re-export multiple items to create a convenient API:

\begin{lstlisting}
// ... (example from std/async/mod.novus - illustrative)

// Re-export future types
pub use std::async::future::Future, Poll, Context, Waker
pub use std::async::future::Ready, Pending, ready, pending

// Re-export executor functions
pub use std::async::executor::block_on_sleep, poll_sleep_once

// Re-export sleep utilities
pub use std::async::sleep::Sleep, sleep, sleep_secs, sleep_millis
\end{lstlisting}

\section{Constants and Static Variables}
\label{sec:constants-statics}

\subsection{Constants (\texttt{const})}
\label{subsec:constants}

Constants are compile-time values that are inlined at every use site:

\begin{lstlisting}
const MAX_SPRITES: u32 = 8
const PI: i32 = 3_14159  // Fixed-point representation
pub const VERSION: u32 = 1

// Constants can use const functions
const fn double(x: i32) -> i32 { x * 2 }
const DOUBLED: i32 = double(21)  // 42 at compile time
\end{lstlisting}

Properties of constants:
\begin{itemize}
    \item Must be computable at compile time
    \item No memory address (inlined at use sites)
    \item Always immutable
    \item Preferred for configuration values
\end{itemize}

\subsection{Static Variables (\texttt{static})}
\label{subsec:static-vars}

Static variables have a fixed memory location and live for the entire program:

\begin{lstlisting}
// Immutable static - initialized data (size inferred from elements)
static SINE_TABLE: [i16] = [
    0, 804, 1608, 2410, 3212, 4011, 4808, // ...
]

// Mutable static - use with caution
static var frame_count: u32 = 0
static var initialized: bool = false
\end{lstlisting}

Properties of static variables:
\begin{itemize}
    \item Have a memory address (stored in \texttt{.data} or \texttt{.bss})
    \item Live for the entire program lifetime
    \item Immutable by default (\texttt{static var} for mutable)
    \item Must be initialized with a constant expression
\end{itemize}

\subsection{Mutable Statics}
\label{subsec:mutable-statics}

Mutable static variables are declared with \texttt{static var}:

\begin{lstlisting}
static var g_chip_pool: ChipPool = ChipPool {
    blocks: null,
    count: 0,
}

pub fn get_pool() -> *ChipPool {
    return &var g_chip_pool
}
\end{lstlisting}

\textbf{Warning:} Mutable statics are dangerous in multitasking environments. For thread-safe global state, consider using:
\begin{itemize}
    \item Atomic types for simple counters
    \item Semaphores for complex data
    \item Immutable statics with interior mutability
\end{itemize}

\section{Project Structure}
\label{sec:project-structure}

Novus uses TOML configuration files to define projects and workspaces.

\subsection{Single Project}
\label{subsec:single-project}

A simple project structure:

\begin{lstlisting}[language=bash]
myapp/
  project.toml      # Project configuration
  src/
    main.novus      # Entry point
    renderer.novus  # Additional module
    sound/
      player.novus  # Submodule
      mixer.novus
\end{lstlisting}

The \texttt{project.toml} file:

\begin{lstlisting}[language=bash]
[package]
name = "myapp"
version = "0.1.0"
type = "cli"             # cli, workbench, dual, library, device
entry = "src/main.novus" # Entry point (optional, defaults to main.novus)
description = "My Amiga application"
authors = ["Your Name"]

[build]
target_cpu = "68020"     # 68000, 68020, 68040, 68060
fpu = "auto"             # soft, 68881, 68040, auto
output = "build"
optimization_level = 2

[paths]
src = "src"
\end{lstlisting}

\subsection{Workspaces}
\label{subsec:workspaces}

Workspaces contain multiple related projects:

\begin{lstlisting}[language=bash]
myworkspace/
  workspace.toml      # Workspace configuration
  mylib/
    project.toml
    src/
      lib.novus
  myapp/
    project.toml
    src/
      main.novus
  examples/
    project.toml
    src/
      demo.novus
\end{lstlisting}

The \texttt{workspace.toml} file:

\begin{lstlisting}[language=bash]
[workspace]
name = "myworkspace"
version = "0.1.0"
members = ["mylib", "myapp", "examples"]  # Directory names of member projects

[workspace.build]
target_cpu = "68020"
fpu = "auto"
optimization_level = 2
\end{lstlisting}

The \texttt{members} array lists directory names of projects within the workspace.

\subsection{Project Types}
\label{subsec:project-types}

Novus supports several project types, each with different startup behavior:

\begin{description}
    \item[\texttt{cli}] Command-line application launched from Shell with \texttt{argc}/\texttt{argv} arguments
    \item[\texttt{workbench}] Application launched by double-clicking icons; receives WBStartup message with file arguments
    \item[\texttt{dual}] Application that handles both CLI and Workbench launches
    \item[\texttt{library}] AmigaOS shared library (\texttt{.library}) with ROMTag and vector table
    \item[\texttt{device}] AmigaOS device driver (\texttt{.device}) with BeginIO/AbortIO interface
\end{description}

\textbf{Important:} Workbench applications \textit{must} reply to the WBStartup message using \texttt{ReplyMsg()} when exiting, or Workbench will hang. Use \texttt{Input() == 0} to detect Workbench launches.

\subsection{Creating Projects}
\label{subsec:creating-projects}

Use the \texttt{novusc new} command to create projects:

\begin{lstlisting}[language=bash]
# Create a new workspace
novusc new myworkspace

# Create a CLI project in the workspace
cd myworkspace
novusc new mytool --type cli

# Create a Workbench application
novusc new mygui --type workbench

# Create a library
novusc new mylib --type library
\end{lstlisting}

\subsection{Building Projects}
\label{subsec:building-projects}

Build commands:

\begin{lstlisting}[language=bash]
# Build all projects in workspace
novusc build

# Build a specific project
novusc build mytool

# Build with specific CPU target
novusc build --cpu 68040
\end{lstlisting}

\section{Standard Library Organization}
\label{sec:stdlib-org}

The Novus standard library is organized into focused modules:

\subsection{Core Modules}
\label{subsec:core-modules}

\begin{description}
    \item[\texttt{std::core}] Fundamental types: \texttt{Option}, \texttt{Result}, \texttt{Drop}, \texttt{Clone}, \texttt{Eq}, etc.
    \item[\texttt{std::prelude}] Common re-exports for convenience
\end{description}

\subsection{Data Structures}
\label{subsec:data-structures}

\begin{description}
    \item[\texttt{std::collections::vec}] Dynamic array (\texttt{Vec<T>})
    \item[\texttt{std::collections::hashmap}] Hash map (\texttt{HashMap<K, V>})
    \item[\texttt{std::collections::hashset}] Hash set (\texttt{HashSet<T>})
    \item[\texttt{std::collections::vecdeque}] Double-ended queue
    \item[\texttt{std::collections::ringbuffer}] Fixed-size ring buffer
    \item[\texttt{std::collections::bitset}] Bit set
    \item[\texttt{std::collections::slotmap}] Slot map for entity storage
    \item[\texttt{std::collections::freelist}] Free list allocator
    \item[\texttt{std::collections::smallvec}] Small-buffer-optimized vector
    \item[\texttt{std::collections::intrusive\_list}] Intrusive doubly-linked list
\end{description}

\subsection{Memory Management}
\label{subsec:memory-modules}

\begin{description}
    \item[\texttt{std::memory::block}] Raw memory blocks (\texttt{MemoryBlock})
    \item[\texttt{std::memory::chip\_pool}] Chip memory pool
    \item[\texttt{std::memory::chip\_cache}] Chip memory cache
\end{description}

\subsection{Strings}
\label{subsec:string-modules}

\begin{description}
    \item[\texttt{std::strings::core}] String slice type (\texttt{Str})
    \item[\texttt{std::strings::format}] Formatting and \texttt{Display} trait
\end{description}

\subsection{Error Handling}
\label{subsec:error-modules}

\begin{description}
    \item[\texttt{std::error::errors}] Error types: \texttt{DosError}, \texttt{ExecError}, \texttt{GraphicsError}, etc.
\end{description}

\subsection{I/O and Networking}
\label{subsec:io-modules}

\begin{description}
    \item[\texttt{std::io}] Basic I/O functions
    \item[\texttt{std::net::socket}] BSD socket interface
    \item[\texttt{std::net::tcp}] TCP client/server
    \item[\texttt{std::net::udp}] UDP sockets
    \item[\texttt{std::net::dns}] DNS resolution
\end{description}

\subsection{Concurrency}
\label{subsec:concurrency-modules}

\begin{description}
    \item[\texttt{std::async}] Async/await runtime
    \item[\texttt{std::sync}] Synchronization primitives
    \item[\texttt{std::thread}] Task management
\end{description}

\subsection{Graphics and UI}
\label{subsec:graphics-modules}

\begin{description}
    \item[\texttt{std::graphics::bitmap}] Bitmap handling
    \item[\texttt{std::graphics::sprite}] Hardware sprites
    \item[\texttt{std::graphics::copper}] Copper list generation
    \item[\texttt{std::graphics::blitter}] Blitter operations
    \item[\texttt{std::graphics::intuition}] Intuition window/screen wrappers
    \item[\texttt{std::ui::window}] Window management
    \item[\texttt{std::ui::screen}] Screen management
    \item[\texttt{std::ui::menu}] Menu creation
    \item[\texttt{std::ui::events}] Event handling
\end{description}

\subsection{AmigaOS FFI}
\label{subsec:ffi-modules}

The FFI modules provide direct bindings to AmigaOS libraries:

\begin{description}
    \item[\texttt{std::ffi::exec}] Exec library (memory, tasks, signals)
    \item[\texttt{std::ffi::dos}] DOS library (files, paths, CLI)
    \item[\texttt{std::ffi::graphics}] Graphics library (rendering, text)
    \item[\texttt{std::ffi::intuition}] Intuition library (windows, screens)
    \item[\texttt{std::ffi::layers}] Layers library (layer management)
    \item[\texttt{std::ffi::gadtools}] GadTools library (GUI toolkit)
    \item[\texttt{std::ffi::utility}] Utility library (tags, hooks)
    \item[\texttt{std::ffi::workbench}] Workbench library (icons, AppIcons)
    \item[\texttt{std::ffi::icon}] Icon library (icon handling)
    \item[\texttt{std::ffi::asl}] ASL library (file/font requesters)
    \item[\texttt{std::ffi::diskfont}] Diskfont library (font loading)
    \item[\texttt{std::ffi::commodities}] Commodities library (hotkeys)
    \item[\texttt{std::ffi::iffparse}] IFFParse library (IFF files)
    \item[\texttt{std::ffi::timer\_device}] Timer device (delays, timing)
    \item[\texttt{std::ffi::audio\_device}] Audio device (Paula access)
    \item[\texttt{std::ffi::bsdsocket}] BSD sockets (networking)
    \item[\texttt{std::ffi::amiga\_structs}] AmigaOS structure definitions
    \item[\texttt{std::ffi::amiga\_consts}] AmigaOS constants
\end{description}

\subsection{Other Modules}
\label{subsec:other-modules}

\begin{description}
    \item[\texttt{std::math}] Fixed-point math, trigonometry
    \item[\texttt{std::audio}] Paula audio, ProTracker player
    \item[\texttt{std::test}] Test framework
    \item[\texttt{std::os}] Operating system utilities
\end{description}

\section{The Prelude}
\label{sec:prelude}

The prelude (\texttt{std::prelude}) re-exports commonly used types:

\begin{lstlisting}
// ... (contents of std/prelude.novus - illustrative)
pub use std::core::Option, Result, Drop, From
pub use std::strings::core::Str
pub use std::error::errors::DosError, ExecError, IntuitionError, GraphicsError
\end{lstlisting}

Import the prelude for quick access to common types:

\begin{lstlisting}
from std::prelude import *

pub fn main() -> i32 {
    let result: Result<i32, DosError> = do_something()
    match result {
        Ok(value) => return value,
        Err(_) => return 1
    }
}
\end{lstlisting}

\section{Best Practices}
\label{sec:module-best-practices}

\subsection{Visibility Guidelines}
\label{subsec:visibility-guidelines}

\begin{enumerate}
    \item \textbf{Default to private} --- only make items \texttt{pub} when needed for external API
    \item \textbf{Use \texttt{internal} for shared implementation} --- better than making everything \texttt{pub}
    \item \textbf{Document public APIs} --- anything \texttt{pub} should have doc comments
    \item \textbf{Keep modules focused} --- one module, one responsibility
\end{enumerate}

\subsection{Import Guidelines}
\label{subsec:import-guidelines}

\begin{enumerate}
    \item \textbf{Prefer specific imports} over wildcards for clarity
    \item \textbf{Use wildcards sparingly} --- only for well-known modules like \texttt{amiga\_consts}
    \item \textbf{Group imports logically} --- core types, then collections, then FFI
    \item \textbf{Use aliases} to avoid name conflicts
\end{enumerate}

\subsection{Project Organization}
\label{subsec:project-org}

\begin{enumerate}
    \item \textbf{Use workspaces} for multi-project development
    \item \textbf{Separate library and application code} into different projects
    \item \textbf{Keep tests in a separate project} that depends on the library
    \item \textbf{Use meaningful directory structure} that matches module paths
\end{enumerate}

\subsection{Constants vs Statics}
\label{subsec:const-vs-static}

\begin{center}
\begin{tabular}{lll}
\textbf{Use} & \textbf{When} \\
\hline
\texttt{const} & Configuration values, magic numbers, version strings \\
\texttt{static} & Lookup tables, large data that shouldn't be inlined \\
\texttt{static var} & Global state (use sparingly, prefer safer alternatives) \\
\end{tabular}
\end{center}

\section{Summary}
\label{sec:modules-summary}

This chapter covered the Novus module system:

\begin{itemize}
    \item \textbf{Imports}: \texttt{from module::path import items}
    \item \textbf{Visibility}: Private (default), \texttt{pub}, \texttt{internal}, \texttt{extern}
    \item \textbf{Re-exports}: \texttt{pub use} for API convenience
    \item \textbf{Constants}: \texttt{const} for compile-time values
    \item \textbf{Statics}: \texttt{static} for fixed-location data
    \item \textbf{Projects}: \texttt{project.toml} and \texttt{workspace.toml}
    \item \textbf{Standard library}: Well-organized modules for all needs
\end{itemize}

Good module organization makes code easier to understand, maintain, and reuse. Start with private visibility and only expose what's needed for your public API.

The next chapter covers error handling with \texttt{Result} and \texttt{Option}.
